\subsection{RQ4: Overhead and Practicality}
\label{sec:rq4}

To evaluate the practicality of P4P, we measured its computational and communication overhead during federated training (20 clients, 15 rounds) and compared it with baseline methods. The results, summarized in Table~\ref{tab:training_time_pivoted}, indicate that P4P introduces negligible additional cost while maintaining high efficiency.

\begin{table}[h!]
\centering
\caption{Server-side computation overhead per round (in seconds) ($\alpha = 0.5$)}
\label{tab:training_time_pivoted}
\begin{tabular}{l | c c c c}
\hline
Attack Type & No-Def & Recess & FedMP & P4P \\
\hline
LIE & 210.58 & 242.37 & 243.21 & \textbf{222.37} \\
Sign Flip & 203.36 & 235.76 & 235.41 & \textbf{224.59} \\
Scaling & 205.64 & 235.21 & 235.13 & \textbf{226.17} \\
Gaussian& 211.06 & 212.11 & 215.18 & \textbf{211.45} \\
Label Flip & 212.46 & 222.97 & 221.21 & \textbf{220.52} \\
Backdoor& 219.28 & \textbf{211.80} & 224.03 & 219.28 \\
\hline
\end{tabular}
\end{table}

Across all attack scenarios, P4P’s total training time is nearly identical to the undefended baseline (\textit{No-Def}). The difference ranges from a 1.0\% decrease (i.e., faster) under the LIE attack to only a 0.6\% increase under the Sign Flip attack, with no measurable overhead observed in the Backdoor setting. In comparison, other defenses such as FedMP incur higher costs, showing up to a +2.2\% increase in training time. These results confirm that P4P achieves robustness without sacrificing computational efficiency.

The lightweight design of P4P contributes to its practicality in two key aspects. First, computational overhead remains low since each client’s high-dimensional model update is reduced to a single scalar probe response via cosine similarity. This compact representation enables efficient anomaly detection on one-dimensional scores instead of full gradients. Second, communication overhead is minimal, requiring only the transmission of one additional probe vector from the server and a single scalar value ($r_i^t$) from each client per round. This added cost is negligible relative to standard model updates in FL.

Overall, these quantitative results demonstrate that P4P satisfies RQ4 by providing a scalable, resource-efficient, and practical defense mechanism that integrates seamlessly into existing FL pipelines.